\section{Conclusion}

Throughout this dissertation, we have pursued two complementary themes in numerical relativity. The first is the numerical modeling of black-hole spacetimes in Einstein--Maxwell--dilaton--axion theory, an extension of general relativity that includes electromagnetic, dilaton, and axion fields. Within this setting, we studied the nonlinear stability of charged black holes and investigated how these additional fields can affect gravitational waveforms from binary black-hole mergers. This work provides an example of how numerical relativity can be used to explore strong-field dynamics in theories beyond general relativity.

The second theme is the improvement of numerical relativity methods through the use of compact finite-difference operators. These methods are designed to improve the accuracy and efficiency of simulations of strong-field gravitational systems. In this final chapter, we summarize the main conclusions from each part of the dissertation and discuss directions for future work.

In the first EMDA portion of this dissertation, we developed a method for constructing approximate binary black-hole initial data in Einstein--Maxwell--dilaton--axion theory. This allowed us to evolve head-on collisions of charged black holes across a range of initial spins, charges, and field configurations. These simulations were used to study whether the additional fields present in EMDA theory remain associated with black holes in the fully nonlinear merger regime.

We found that scalar hair persists through merger for initially scalarized Kerr--Sen black holes. When the black holes carry nontrivial electric charge, the dilaton field remains present after merger. When both charge and spin are present, the remnant also retains an axion field. The horizon quantities of the final black hole settle to approximately constant values, supporting the interpretation that the remnant approaches another Kerr--Sen black hole. We also showed that initially unscalarized black holes can dynamically develop scalar hair. In particular, initially Kerr--Newman-like black holes in EMDA evolve toward configurations with nontrivial scalar structure. These results suggest that scalar and axionic hair can be robust features of charged black holes in EMDA theory, rather than artifacts of the initial data.

For inspiraling black-hole binaries, we found that increasing the charge produces small changes in the gravitational waveform relative to the corresponding neutral GR case. These waveform differences provide a first indication of how EMDA fields may affect binary black-hole signals in the strong-field regime. We also constructed a modest waveform catalog from the simulations performed in this dissertation, including several charge configurations, spins, and mass ratios. In particular, this catalog includes EMDA binary black-hole mergers with mass ratios as large as $q=8$, demonstrating that these simulations can be extended beyond the equal-mass case.

In the broader EMDA theory-space study, we investigated how binary black-hole evolutions change as the coupling parameters are varied. We explored a range of EMDA coupling choices and compared these results with simulations in Einstein--Maxwell theory, Kaluza--Klein-inspired theory, and related Einstein--Maxwell--dilaton models. This allowed us to search for significant differences in the dynamics and to assess how theory-dependent effects appear in numerical evolutions. For the small charges considered in that study, the dominant gravitational-wave signal was only weakly affected by the coupling choice, with the largest differences appearing in the strongest-coupling cases.

In the compact finite-difference work, we developed novel compact derivative operators and evaluated their performance against comparable explicit finite-difference operators. In the tests considered here, the fourth-order compact operators achieved accuracy comparable to fourth-order explicit operators while using roughly half the resolution. We also demonstrated that sixth-order compact derivatives can outperform their explicit counterparts. These improvements were shown for both first- and second-derivative operators and were tested across three systems: Maxwell theory, a nonlinear sigma model, and the BSSN equations. The improved accuracy of these operators suggests that compact finite differences may provide a useful tool for improving the fidelity of numerical relativity simulations.

Several directions remain for future work. First, compact finite-difference operators should be implemented and tested in full binary black-hole evolutions. Preliminary work in this direction has already shown some advantages, but additional simulations are needed to determine whether these operators provide a robust improvement over standard explicit finite differences in realistic merger calculations. If these advantages persist, compact finite differences could provide a path toward faster and higher-fidelity gravitational-wave simulations, which will be increasingly important as next-generation detectors require more accurate waveform models.

Second, future EMDA simulations could be used to make closer contact with gravitational-wave observations. One promising approach would be to simulate systems similar to high-significance observed binary black-hole mergers whose masses and spins are relatively well constrained. By comparing EMDA waveforms against observed signals, one could place constraints on the possible dilaton and axion charges of the merging black holes. Such a study would help connect the theoretical and numerical results of this dissertation to observational tests of deviations from general relativity.

Taken together, the results of this dissertation demonstrate how numerical relativity can both expand the study of black-hole dynamics beyond general relativity and improve the computational tools needed to model gravitational-wave sources with increasing accuracy.
